\documentclass[aspectratio=169]{beamer}
\usepackage[utf8]{inputenc}
\usepackage[T1]{fontenc}
\usepackage[spanish]{babel}
\usepackage{tikz}
\usetheme{Madrid}
\definecolor{navy}{RGB}{19,42,65}
\definecolor{teal}{RGB}{0,116,123}
\setbeamercolor{structure}{fg=teal}
\setbeamercolor{palette primary}{bg=navy,fg=white}
\setbeamertemplate{navigation symbols}{}
\title[Proyecto 10 / DSM]{Memoria compartida distribuida\\Mini-simulador a nivel de página}
\subtitle{Dos procesos, páginas que migran y escrituras exclusivas}
\author[Equipo 10]{Armando, Daniel, Guillermo, Daniel y Mauricio}
\institute[SO]{Sistemas Operativos\\Catedrático: David Garcia Cerna\\Institución y sección: por completar}
\date{17 de octubre de 2026}
\begin{document}
\begin{frame}\titlepage\end{frame}

\begin{frame}{El problema que resolvemos}
\begin{itemize}
\item Cada proceso tiene su propia memoria: no ve directamente el arreglo del otro.
\item Queremos consultar y modificar páginas como parte de un espacio lógico común.
\item Los datos viajan por sockets; las reglas deciden quién puede usarlos.
\end{itemize}
\begin{block}{Objetivo de la guía}
Demostrar migración para lectura e invalidación antes de escribir.
\end{block}
\end{frame}

\begin{frame}{Qué construimos}
\begin{columns}
\column{.48\textwidth}
\begin{itemize}
\item C11 en Pop!\_OS / Linux.
\item Dos procesos independientes.
\item TCP IPv4: 127.0.0.1.
\item Un hilo por proceso; poll.
\end{itemize}
\column{.48\textwidth}
\begin{itemize}
\item 1 MiB local por nodo.
\item 256 páginas de 4096 bytes.
\item Inicio: N0 dueño, N1 inválido.
\item Fallo simulado al revisar etiquetas.
\end{itemize}
\end{columns}
\begin{alertblock}{Alcance}
No se interceptan punteros C ni fallos reales del kernel.
\end{alertblock}
\end{frame}

\begin{frame}{Una conexión y un coordinador}
\centering
\begin{tikzpicture}[every node/.style={align=center}]
\node[draw=teal,thick,rounded corners,minimum width=4cm,minimum height=2cm] (n0) at (0,0)
{N0: coordinador\\directorio + memoria\\puerto 5000};
\node[draw=teal,thick,rounded corners,minimum width=4cm,minimum height=2cm] (n1) at (8,0)
{N1: participante\\memoria local\\puerto 5001};
\draw[<->,thick] (n0) -- node[above]{TCP: control y páginas} (n1);
\end{tikzpicture}
\vspace{.4cm}
\begin{itemize}
\item N0 conecta; N1 acepta. Ambos pueden iniciar primero.
\item El coordinador permanece en N0; la propiedad de datos puede migrar.
\end{itemize}
\end{frame}

\begin{frame}{Tres etiquetas por página}
\begin{description}
\item[INVALID] Puede haber bytes viejos, pero no se permite leerlos.
\item[READ\_ONLY] Copia válida para lectura; no autoriza modificar.
\item[OWNER] Copia autoritativa; necesita exclusividad antes de escribir.
\end{description}
\begin{block}{Dos lecturas distintas}
\texttt{read}: migra y deja un dueño único.\\
\texttt{read-copy}: conserva una réplica para mostrar invalidación de lectores.
\end{block}
\end{frame}

\begin{frame}{Demo 1: la página migra para leer}
\begin{enumerate}
\item N0: \texttt{write 3 hola desde N0}.
\item N1: \texttt{read 3}; detecta INVALID y solicita la página.
\item Se preservan datos; N0 invalida su copia.
\item Viajan 4096 bytes; N1 instala, lee y confirma.
\end{enumerate}
\begin{block}{Comprobar en ambas terminales}
\texttt{status 3}: N0 = INVALID; N1 = OWNER; versión = 1.
\end{block}
\end{frame}

\begin{frame}{Demo 2: invalidar antes de escribir}
\begin{enumerate}
\item N0: \texttt{read-copy 3}; queda READ\_ONLY y N1 conserva OWNER.
\item N1: \texttt{write 3 actualizado por N1}; envía INVALIDATE.
\item N0 aplica INVALIDATE\_APPLIED y concede exclusividad.
\item Solo entonces N1 escribe; versión = 2.
\item N0: \texttt{read 3}; recupera el texto nuevo y la propiedad.
\end{enumerate}
\begin{block}{Otra dirección}
Si N0 escribe con lector N1, envía INVALIDATE y espera ACK remoto antes de escribir.
\end{block}
\end{frame}

\begin{frame}{Por qué no se usan datos viejos}
\begin{itemize}
\item Toda operación pasa por N0, incluso un acierto local.
\item Se admite una transacción global: obtener datos, invalidar, conceder, confirmar.
\item La copia INVALID nunca sirve una lectura.
\item La escritura requiere confirmaciones; cada escritura aumenta la versión.
\item Otra solicitud recibe BUSY y se reintenta manualmente.
\end{itemize}
\begin{block}{Distinción}
Coherencia trata las copias; consistencia trata el orden observable de accesos.
\end{block}
\end{frame}

\begin{frame}{Transporte y errores}
\begin{itemize}
\item TCP entrega bytes, no mensajes completos.
\item Cabecera de 40 bytes y longitudes exactas; una página lleva 4096 bytes.
\item Buffers conservan recepción y envío parciales; no bloquean la terminal.
\item Se rechazan índices, tamaños, versiones o ACK fuera de fase.
\item Plazos: saludo/config 10 s; operación/trama/cola 5 s.
\end{itemize}
\begin{alertblock}{Ante caída}
Se aborta y se reinician ambos nodos. No se promete recuperación de propiedad.
\end{alertblock}
\end{frame}

\begin{frame}{Pruebas y resultados reales}
\begin{itemize}
\item Compilación limpia sin warnings; tres programas de pruebas C.
\item 22 pruebas de integración aprobadas con procesos TCP reales.
\item Nueve escenarios DSM con ASan/UBSan, sin errores detectados.
\item ACK retrasado: se comprueba que la escritura espere y que otra solicitud reciba BUSY.
\item 48 operaciones con semilla fija contrastadas con un modelo de valores.
\end{itemize}
\begin{block}{Demo de cinco operaciones}
3 transferencias de página = 12288 bytes de payload; 3 fallos simulados;
3 invalidaciones aplicadas; versión final 2; propietario final N0.
\end{block}
\end{frame}

\begin{frame}[fragile]{Reproducir desde la terminal}
\begin{verbatim}
make
# Terminal N0
./build/dsm-node --id 0 --port 5000 --peer-port 5001
# Terminal N1
./build/dsm-node --id 1 --port 5001 --peer-port 5000
\end{verbatim}
Esperar DSM\_READY. Introducir una orden y esperar su resultado.
\begin{verbatim}
make test      # pruebas
make demo      # demo y trazas verificadas
make package   # fuentes + articulo y presentacion PDF
\end{verbatim}
\end{frame}

\begin{frame}{Límites, aprendizaje y fuentes}
\begin{itemize}
\item Dos nodos, memoria volátil, loopback y una operación a la vez.
\item Sin acceso transparente, persistencia, reconexión ni tolerancia a fallos.
\item Las pruebas dan evidencia; no son una demostración formal exhaustiva.
\item Aprendizaje: relacionar páginas, propiedad, mensajes y orden de acceso.
\end{itemize}
\small Referencias: Li y Hudak, ACM TOCS 7(4), 1989; Lamport, IEEE TC C-28(9), 1979;
Linux man-pages: tcp(7), recv(2), send(2), poll(2). Enlaces en el artículo y FUENTES.md.
\end{frame}
\end{document}
